\chapter{Background and Literature Review}
\label{chap:background}

This chapter reviews the relevant literature across four areas. First, checkpoint formats and how their storage layout affects loading behaviour. Second, serverless and multi-tenant inference systems and the cold-start problem. Third, the Linux I/O interfaces available to applications, with particular attention to io\_uring. Fourth, NVMe storage behaviour and the role of memory-mapping in checkpoint loading. For each area the review identifies what prior work assumes, what it does not examine, and how this project differs from it.

\section{Checkpoint Formats and Loading Behaviour}

LLM weights are stored on disk in a structured format. The format determines how loading software reads tensor data: whether reads are sequential or scattered, whether they require many file operations or few, and whether the layout helps or hinders the storage device in maintaining high queue depth.

SafeTensors is a format introduced by Hugging Face for safe model distribution~\cite{safetensors_docs}. The file begins with a JSON header describing each tensor, including its name, dtype, shape, and byte offset. All tensor data follows the header contiguously. Once the header is parsed, loading software can read each tensor by seeking to its declared offset and reading the declared number of bytes. The format has two properties that are relevant here. First, it avoids loading the entire file into memory before parsing, unlike Python pickle. Second, it prohibits unsafe deserialisation, which matters when loading untrusted checkpoints. From a systems perspective, however, SafeTensors leaves one important question open: it standardises the layout, but not the loading strategy. The format tells the loader where the bytes are, but not how aggressively it should submit reads to the operating system.

PyTorch checkpoint files typically use Python pickle serialisation. A model is saved by calling \texttt{torch.save()}, which serialises both the model state and the pickle protocol metadata. This has well-known safety problems: a malicious checkpoint can execute arbitrary code on load~\cite{safetensors_docs}. It also requires reading the entire file before unpacking, which is impractical for large models where host RAM may not hold the full checkpoint alongside the deserialised tensors. These limitations are not just theoretical: several public model repositories have had to clean up malicious checkpoints. SafeTensors was designed partly to address both concerns.

GGUF is a more recent format designed for quantised models and memory-mapped loading. It stores tensors in a layout optimised for direct memory mapping and supports multiple quantisation schemes. GGUF tensors are organised so that the memory layout of the file mirrors the layout the model expects at runtime, reducing or eliminating a copy step. This is a different trade-off from SafeTensors: GGUF prioritises memory-mapping friendliness over simple sequential loading.

Partitioned checkpoint formats distribute tensor data across multiple files. ServerlessLLM uses this approach to support partial model loading: only the files needed for the current request are loaded, rather than the entire checkpoint~\cite{FU:2024,serverlessllm_repo}. A central index maps each tensor name to its partition file and byte offset. Loading a complete checkpoint requires opening each partition file and reading from the specified offset. The advantage is that loading can be parallelised across files and across network boundaries if partitions are distributed. The disadvantage is that the access pattern becomes more fragmented, with per-file overhead accumulating. This makes partitioned formats especially interesting from a backend-selection perspective, because fragmented range reads stress the submission path differently from large contiguous shard reads.

Table~\ref{tab:format_comparison} summarises how these formats differ in ways relevant to loading speed.

\begin{table}[h]
\caption{Comparison of checkpoint formats relevant to loading performance.}
\label{tab:format_comparison}
\centering
\begin{tabular}{p{3.5cm}p{3cm}p{3cm}p{3.5cm}}
\toprule
Format & Layout & Loading behaviour & Safety properties \\
\midrule
PyTorch pickle & Single file, full load & Entire file read then unpacked & Unsafe: arbitrary code execution possible \\
SafeTensors & Single file, header then data & Header parsed, individual tensor reads at offset & Safe: no deserialisation of executable code \\
GGUF & Single or split file, memory-map oriented & mmap-friendly layout, minimal copy & Safe; quantisation-aware \\
ServerlessLLM partitioned & Multiple partition files, central index & Per-partition open and read at offset & Safe; supports partial loading \\
\bottomrule
\end{tabular}
\end{table}

What these formats share is that loading is fundamentally a series of read operations with known byte ranges. The storage system sees a set of read requests. Whether those requests are processed efficiently depends on how they are submitted.

The existing literature describes format design well but focuses on safety, portability, and memory layout. It does not systematically examine how the I/O submission path affects the effective loading speed of checkpoints that are already in a safe format. This report addresses that gap by keeping the checkpoint abstraction fixed while varying the execution policy used to fetch the underlying bytes.

\section{Serverless and Multi-Tenant Inference Systems}

Serverless inference platforms such as AWS Lambda, Google Cloud Functions, and specialised services run LLM inference on demand. Memory and GPU resources are allocated when a request arrives and freed when the request completes. Between requests, resources may be handed to other tenants. When a model is not resident in memory, the platform must load it from persistent storage before inference can begin. This delay is the cold-start latency.

Cold-start latency in serverless settings has been studied from several angles. One class of solutions addresses it through scheduling and placement: keeping models warm, pre-loading them on a pool of standby instances, or migrating running instances to avoid a load altogether~\cite{FU:2024,serverlessllm_repo}. Another class addresses it through model optimisation: quantisation, pruning, and early exit strategies that reduce the amount of data that must be loaded or processed. A third class, which this project falls into, examines whether the loading mechanism itself can be made faster.

Snapshot-based warm-up techniques preload a model into memory before the first request arrives~\cite{silva2021snapshot}. Prebaking keeps a function instance warm by periodically sending requests to it. SEUSS optimises the warm-up process for serverless function containers. These approaches treat the checkpoint loading step as a fixed cost and work around it rather than reducing it.

The cold-start problem becomes more acute as model sizes grow. A seven-billion-parameter model in float16 requires roughly fourteen gigabytes of storage and a similar amount of RAM or VRAM. Loading this much data from an NVMe SSD takes hundreds of milliseconds to several seconds depending on the device and submission path. In a production serving system, this is a large fraction of the total latency budget, particularly for interactive applications.

vLLM is a high-throughput inference engine that supports paged memory management and continuous batching~\cite{pagedattention,vllm_repo}. Its design is representative of a broader class of serving systems, including ORCA~\cite{orca}, where throughput improvements come from scheduling, batching, and runtime memory management rather than from a specialised checkpoint-loading path alone. This makes vLLM a useful integration target in this project: if backend-level gains matter in practice, they should survive at least partially into a serving engine with its own substantial initialisation overhead.

Most cold-start research focuses on what happens before loading begins or after it ends. Loading speed itself is often treated as a constant or folded into a broader serving metric. This report instead treats loading as a first-class systems problem and asks which backend policy best matches a given checkpoint regime.

\section{Linux I/O Interfaces}

Applications running on Linux have several interfaces for reading from storage.

The most common is POSIX synchronous I/O, based on the \texttt{read()} system call. A thread calls \texttt{read(fd, buf, n)} and blocks until the data is available in the buffer. To issue multiple reads concurrently, an application must use multiple threads. Each thread is blocked while its read is in flight. Thread creation and context switching have costs, so thread pool sizes are limited. This limits the number of outstanding reads that a synchronous application can maintain at once.

Linux AIO was introduced to provide genuine asynchronous I/O without threads. However, Linux AIO has significant practical limitations. It does not support most filesystems including ext4, XFS, and btrfs. It requires manual error handling at a low level. The API is complex and brittle to use correctly~\cite{MAN7_AIO7}. Few production applications use it directly for these reasons.

POSIX AIO is a userspace library implementation of asynchronous I/O using a thread pool under the hood. It inherits the limitations of threads and adds implementation complexity. It is not widely used in new systems code.

io\_uring was introduced in Linux kernel 5.1 and designed to address the limitations of earlier async I/O interfaces~\cite{AXBOE_IOURING_PDF,MAN7_IOURING7}. Instead of system calls for each operation, applications interact with io\_uring through two shared-memory rings: a submission queue and a completion queue. Submission queue entries (SQEs) are prepared in memory and submitted with a single \texttt{io\_uring\_enter} system call, or in kernel-polling mode with no system call at all. Completion queue entries (CQEs) are retrieved by polling the completion queue.

io\_uring provides several features relevant to checkpoint loading. Registered buffers are pre-allocated and registered with the kernel, eliminating per-request memory allocation overhead. Fixed buffers also reduce the number of data copies per operation. Batched submission allows many requests to be prepared and submitted together, reducing per-request syscall overhead. Polled completion allows an application to busy-wait on completions without system call overhead when latency is critical. The low-level Rust ecosystem mirrors this design space: Tokio provides a mature async runtime~\cite{tokio_docs,tokio_repo}, while the \texttt{io-uring} crate~\cite{io_uring_crate} exposes the lower-level ring interface needed to experiment with more explicit submission strategies.

Recent work has examined when io\_uring helps in database systems. Jasny et al. studied io\_uring integration into a buffer manager and a network-shuffling workload, finding that io\_uring benefits depend heavily on workload characteristics and queue depth~\cite{IOURING_DATABASE_GUIDELINES}. Their results show that io\_uring helps most when there are many small to medium-sized requests and the device queue depth is high. For bulk sequential reads, the advantage is smaller because the syscall overhead is already low relative to the data transfer time. They note that naive replacement of synchronous I/O with io\_uring does not automatically yield improvements; the workload must exercise the features that differentiate io\_uring. This observation is one of the key literature anchors for the present report: backend family alone is not enough, and execution strategy must be matched to workload shape.

This observation is directly relevant to checkpoint loading. A contiguous SafeTensors file with large sequential reads may not benefit as much from io\_uring as a fragmented partitioned layout with many small reads. The experiment section of this report examines this question directly, and the final results confirm that backend selection must be heuristic and workload-aware rather than fixed.

The io\_uring interface itself has limitations relevant to checkpoint loading. Registered buffers work best when all buffers are the same size, but checkpoint tensors vary. Managing variable-sized buffers with fixed registration requires splitting tensors into fixed-size chunks, which adds complexity. The kernel-polling mode consumes CPU even when the device is idle. The interrupt-driven mode has lower CPU overhead but higher latency per completion. There is no single optimal mode for all checkpoint layouts.

\section{NVMe Storage and Queueing Behaviour}

Modern NVMe devices expose multiple submission and completion queues directly to the host~\cite{NVME_SPEC}. The device processes commands from these queues concurrently using internal parallelism through parallel flash die access and controller-side scheduling. To exploit this parallelism, the host must submit many commands simultaneously. Devices are optimised for sustained throughput under high queue depth; under low queue depth they frequently sit idle between completions.

The Linux block layer manages I/O requests before they reach the NVMe driver. The blk-mq (multiqueue block layer) maps I/O requests to the device's hardware submission queues based on block offset. This mapping is designed to keep related requests on the same hardware queue and exploit device-level parallelism.

Storage queueing theory explains why throughput under load depends on queue depth. The dominant factor is often not the raw bandwidth of the storage medium but the ability of the host software stack to keep the device's internal queues populated~\cite{STORAGE_QUEUEING}. If the submission path introduces idle gaps, the device cannot compensate even if its peak throughput is high. This is the theoretical basis for the submission bottleneck hypothesis.

\section{Memory-Mapped I/O and the Page Cache}

Linux supports memory-mapped file access through the \texttt{mmap} system call. When a file is mapped, its contents appear in the process address space and can be read like ordinary memory. The operating system loads pages from disk on demand through page faults, rather than in response to explicit \texttt{read()} calls.

For checkpoint loading, mmap has two apparent advantages. First, the kernel handles loading pages in response to faults, so the application does not explicitly submit individual reads. Second, the page cache retains recently accessed checkpoint pages, reducing subsequent load times for frequently accessed models.

However, mmap introduces its own costs. Page faults have overhead, particularly for non-resident pages that must be fetched from disk. For a large checkpoint with many tensors, the cumulative fault overhead can be significant. The kernel's readahead mechanism attempts to prefetch pages it predicts will be accessed next, but it relies on sequential access patterns. Fragmented or random-access layouts may confuse readahead, causing it to load pages that are not used before they are evicted.

The page cache complicates benchmarking in a specific way. If a checkpoint is already in the page cache, a subsequent load will read from RAM rather than from disk, producing misleadingly fast results. Rigorous checkpoint loading evaluation must control for page cache state by evicting relevant pages before each cold measurement run. Without this control, benchmarks measure cache hit rates rather than storage performance.

Direct I/O bypasses the page cache entirely, sending reads directly to the storage device. This is useful for applications that manage their own caching and want predictable behaviour. However, direct I/O requires aligned buffers and is incompatible with memory mapping. For checkpoint loading, direct I/O combined with explicit read submission gives the application most control over the submission path.

\section{Positioning of This Work}

The literature does not yet examine checkpoint loading specifically as an adaptive I/O selection problem. Format design is well understood but focuses on safety and memory layout rather than I/O submission behaviour. Cold-start solutions focus on scheduling and placement rather than loading speed. io\_uring has been studied in databases and file servers but not in the context of LLM checkpoint loading. Memory-mapping trade-offs are understood in general systems terms but not evaluated against synchronous and async I/O in a controlled checkpoint-loading workload.

This report occupies the intersection of these areas. It isolates the I/O submission problem by holding the format and request trace constant while varying the backend mechanism, but its final contribution is broader than a simple backend race. It evaluates backends under cold-start conditions where the page cache does not hold checkpoint data. It examines both contiguous and fragmented checkpoint layouts to understand how format design interacts with submission efficiency. And it integrates the loading framework with a real serving engine to understand how submission gains propagate through a full inference stack. The resulting contribution is a practical adaptive heuristic for checkpoint loading rather than a single universal backend recommendation.

The work contributes a controlled experimental basis for understanding whether the I/O submission path is a meaningful bottleneck in LLM checkpoint loading, under what conditions it matters most, and what design choices in checkpoint formats and serving infrastructure interact with it.
